%*******************************************************
% Abstract in German
%*******************************************************
\begin{otherlanguage}{ngerman}
	\pdfbookmark[0]{Zusammenfassung}{Zusammenfassung}
	\chapter*{Zusammenfassung}
	\begin{sloppypar}
	\noindent \acs{APT}-Scanner müssen kompromittierte oder 
	verdächtige Systeme möglichst schnell analysieren können. 
	Im täglichen Einsatz außerhalb akuter Sicherheitsvorfälle 
	sollen sie neben einer kurzen Scanzeit auch 
	eine geringe Ressourcenauslastung gewährleisten. 
	\vspace{7.56mm}

	\noindent Diese Arbeit untersucht die Performance-Optimierung eines \acs{APT}-Scanners durch die Umstellung von einem objektbasierten, rekursiv-synchronen Scheduling auf ein zentrales, taskbasiertes Scheduling. 

	\vspace{7.56mm}

	\noindent Implementiert und evaluiert wurden drei Scheduling-Algorithmen: First-In, First-Out (\acs{FIFO}), Shortest-Subtree-Time-First (\acs{SSTF}) und Work-Stealing mit \acs{SSTF}-Priorisierung. 
	Die Bewertung erfolgte experimentell in einer kontrollierten Testumgebung auf Debian~13 und Windows~11 mit einem, vier und zehn Threads sowie 20 Wiederholungen je Konfiguration. 
	Gemessen wurden Laufzeit, \acs{CPU}- und \acs{RAM}-Utilization, Waiting Time, Burst Time und der Anteil synchron ausgeführter Tasks. 

	\vspace{7.56mm}

	\noindent Unter Windows zeigen die neuen Scheduler bei vier und zehn Threads klare Laufzeitgewinne gegenüber der Referenz. 
	Bei vier Threads betragen diese Laufzeitgewinne je nach Verfahren 5,69\,\% bis 8,39\,\% und bei zehn Threads 8,89\,\% bis 10,38\,\%. 
	Unter Linux lässt sich dagegen kein konsistenter Laufzeitvorteil nachweisen. 
	\acs{SSTF} liefert die stabilste Gesamtbalance aus Laufzeit, Waiting Time und Speichereigenschaften.

	\vspace{7.56mm}

	\noindent Die Ergebnisse belegen, dass der Wechsel zu einem zentralen, taskbasierten Scheduling möglich ist. 
	Reproduzierbare Leistungsgewinne, besonders unter hohen Lastprofilen, werden ermöglicht.
	\end{sloppypar}

\end{otherlanguage}
